\documentclass[12pt,a4paper]{report}

% ── Essential Packages ──────────────────────────────────────────────────────────
\usepackage[a4paper, margin=1in, top=1in, bottom=1in, left=1in, right=1in]{geometry}
\usepackage{setspace}
\onehalfspacing % Required 1.5 line spacing

\usepackage{mathptmx} % Times New Roman font
\usepackage[utf8]{inputenc}
\usepackage[T1]{fontenc}
\usepackage{amsmath,amssymb,amsfonts}
\usepackage{graphicx}
\usepackage{booktabs}
\usepackage{array}
\usepackage{tabularx}
\usepackage{multirow}
\usepackage{hyperref}
\usepackage{cite}
\usepackage{fancyhdr}
\usepackage{titlesec}
\usepackage{caption}
\usepackage{subcaption}
\usepackage{listings}
\usepackage{xcolor}
\usepackage{microtype}

% ── Hyperlink Styling ──────────────────────────────────────────────────────────
\hypersetup{
    colorlinks=true,
    linkcolor=black,
    citecolor=blue,
    urlcolor=blue,
    pdftitle={Deep Learning-Based Surface Defect Detection - AIC2022 Report},
    pdfauthor={Mo Sahil}
}

% ── Code Listing Styling ───────────────────────────────────────────────────────
\definecolor{codegray}{rgb}{0.5,0.5,0.5}
\definecolor{codepurple}{rgb}{0.58,0,0.82}
\definecolor{backcolour}{rgb}{0.97,0.97,0.98}

\lstdefinestyle{mystyle}{
    backgroundcolor=\color{backcolour},   
    commentstyle=\color{codegray},
    keywordstyle=\color{magenta},
    numberstyle=\tiny\color{codegray},
    stringstyle=\color{codepurple},
    basicstyle=\ttfamily\footnotesize,
    breakatwhitespace=false,         
    breaklines=true,                 
    captionpos=b,                    
    keepspaces=true,                 
    numbers=left,                    
    numbersep=5pt,                  
    showspaces=false,                
    showstringspaces=false,
    showtabs=false,                  
    tabsize=2
}
\lstset{style=mystyle}

% ── Header & Footer Setup ──────────────────────────────────────────────────────
\pagestyle{fancy}
\fancyhf{}
\fancyhead[L]{\small\textit{AIC2022: Introduction to Artificial Intelligence}}
\fancyhead[R]{\small\textit{Course Assignment Report}}
\fancyfoot[C]{\thepage}
\renewcommand{\headrulewidth}{0.4pt}
\renewcommand{\footrulewidth}{0.4pt}

% Fix chapter title spacing
\titleformat{\chapter}[display]
  {\normalfont\bfseries\LARGE}{\chaptertitlename\ \thechapter}{14pt}{\LARGE}
\titlespacing*{\chapter}{0pt}{-20pt}{20pt}

% ──────────────────────────────────────────────────────────────────────────────
\begin{document}

% ══════════════════════════════════════════════════════════════════════════════
% COVER PAGE
% ══════════════════════════════════════════════════════════════════════════════
\begin{titlepage}
    \centering
    \vspace*{0.5cm}
    
    {\Large \textbf{ZAKIR HUSAIN COLLEGE OF ENGINEERING \& TECHNOLOGY}}\\[0.25cm]
    {\large \textbf{ALIGARH MUSLIM UNIVERSITY, ALIGARH}}\\[0.2cm]
    {\normalsize \textbf{Interdisciplinary Center for Artificial Intelligence (ICAI)}}\\[0.6cm]
    
    \hrule height 1.5pt
    \vspace{0.8cm}
    
    {\LARGE \textbf{Deep Learning-Based Surface Defect Detection for Industrial Quality Control}}\\[0.6cm]
    
    {\large \textbf{A Course Assignment Report on Intelligent Software Development}}\\[0.4cm]
    {\normalsize Submitted in partial fulfilment of the requirements for the course:}\\[0.2cm]
    {\large \textbf{Introduction to Artificial Intelligence (AIC2022)}}\\[0.2cm]
    {\normalsize Academic Year: 2026--2027}\\[1.0cm]
    
    \begin{table}[h]
    \centering
    \large
    \begin{tabular}{ll}
        \textbf{Student Name:}    & Mo Sahil \\[0.25cm]
        \textbf{Faculty Number:}  & 25AIBEA101 \\[0.25cm]
        \textbf{Enrollment No:}   & GP0884 \\[0.25cm]
        \textbf{Course Code:}     & AIC2022 \\[0.25cm]
        \textbf{Assignment:}      & Intelligent Software Development \\[0.25cm]
        \textbf{Domain:}          & Computer Vision (Defect Inspection) \\[0.25cm]
        \textbf{Course Instructor:} & \textbf{Dr. Usmani Ahmad Usmani} \\[0.25cm]
        \textbf{Submission Date:} & November 20, 2026 \\
    \end{tabular}
    \end{table}
    
    \vfill
    \hrule height 0.8pt
    \vspace{0.4cm}
    {\normalsize Department of Computer Engineering / ICAI\\Aligarh Muslim University, Aligarh, Uttar Pradesh, India}
\end{titlepage}

% ══════════════════════════════════════════════════════════════════════════════
% ABSTRACT & PRELIMINARIES
% ══════════════════════════════════════════════════════════════════════════════
\pagenumbering{roman}
\setcounter{page}{1}

\begin{abstract}
\onehalfspacing
Automated surface defect inspection is a vital operational challenge in high-throughput discrete manufacturing. Traditional manual visual inspection relies heavily on human visual vigilance, which degrades rapidly due to physiological fatigue, cognitive load, and environmental illumination fluctuations. In this project, an end-to-end intelligent software application is designed, implemented, and empirically evaluated for automated casting product surface defect detection. The proposed system employs transfer learning utilizing a deep Residual Network (ResNet-18) backbone pretrained on the ImageNet benchmark. High-level semantic features are fine-tuned via selective unfreezing of the terminal residual block and custom binary classification head, accompanied by comprehensive spatial and photometric data augmentations. Evaluated on an empirical benchmark of 715 unseen industrial pump impeller casting specimens from a manufacturing production environment, the system achieved a test accuracy of \textbf{99.58\%}, a defect precision of \textbf{100.00\%} (zero false alarms on non-defective components), a defect recall of \textbf{99.34\%}, and an overall F1-score of \textbf{0.9958}, converging in only 5 training epochs. The trained intelligence is packaged and deployed as a real-time web interface built with Gradio, providing operators with single-click visual predictions and calibrated posterior class probabilities.
\end{abstract}

\vspace{0.5cm}
\noindent\textbf{Keywords:} Artificial Intelligence, Computer Vision, Surface Defect Detection, Transfer Learning, ResNet-18, Convolutional Neural Networks, Industrial Quality Control.

\newpage
\tableofcontents
\newpage
\listoffigures
\vspace{0.8cm}
\listoftables
\newpage

\pagenumbering{arabic}
\setcounter{page}{1}

% ══════════════════════════════════════════════════════════════════════════════
% CHAPTER 1: INTRODUCTION
% ══════════════════════════════════════════════════════════════════════════════
\chapter{Introduction}

\section{Background}
Manufacturing quality control represents a foundational pillar of modern industrial engineering and operations management. Across vital economic sectors including automotive powertrains, aerospace turbomachinery, consumer appliances, and heavy structural castings, the mechanical integrity of individual cast components directly dictates system reliability, assembly safety, and operational longevity. Undetected material discontinuities—such as shrinkage porosity, blowholes, slag inclusions, and micro-cracks—inevitably compromise tensile strength, fatigue endurance, and pressure containment limits. Empirical industry audits consistently indicate that defective parts escaping factory inspection gates account for 5\% to 15\% of total manufacturing costs via warranty remediation, product recalls, unplanned assembly stoppages, and lasting reputational damage \cite{stamatis2003failure}.

For decades, industrial defect classification has depended upon manual human visual examination performed at production-line inspection stations. Although human visual perception possesses impressive contextual flexibility, it is fundamentally constrained by well-documented psychological and ergonomic limitations. Human inspection accuracy degrades steeply under repetitive visual stimuli; cognitive psychology literature establishes that human inspector error rates escalate dramatically after two to three hours of continuous line monitoring \cite{drury2012inspection}. Furthermore, manual inspection throughput is strictly bounded by human cognitive cycle times, creating severe bottlenecks in high-speed automated production cells.

Early efforts to automate surface inspection utilized classical machine vision techniques predicated upon handcrafted feature representations. Operators calibrated mathematical spatial kernels, including Sobel and Canny gradient filters \cite{canny1986computational}, morphological structuring elements, Gabor wavelets, and Haralick Gray-Level Co-occurrence Matrices (GLCM) \cite{haralick1973textural}. While computationally modest, handcrafted vision algorithms exhibit severe fragility. Subtle variations in surface reflectivity, ambient lighting angle, part positioning, or cast metal texture routinely provoke catastrophic false-positive spikes or missed detections. Whenever a casting geometry changes, classical filter parameters require laborious manual re-engineering.

\section{Motivation}
The advent of deep Convolutional Neural Networks (CNNs) has revolutionized computer vision by learning hierarchical, data-driven feature extractors directly from raw pixel matrices \cite{krizhevsky2012imagenet, simonyan2015very}. Modern Industry 4.0 frameworks require automated inspection architectures that deliver high classification accuracy, robust generalisation across lighting variations, and low latency suitable for real-time edge deployment.

However, training deep neural networks from random weight initializations typically necessitates hundreds of thousands of annotated images to avoid severe overfitting. In discrete manufacturing, acquiring massive curated libraries of defective components is economically prohibitive and practically difficult because defect events are comparatively rare. Transfer learning addresses this asymmetry by reusing generalized feature extractors learned from vast natural image corpora (such as ImageNet \cite{russakovsky2015imagenet}) and adapting high-level representations to industrial target domains with small-to-medium labeled datasets \cite{oquab2014learning}. This strategy provides the primary technological motivation for the present project.

\section{Scope of the Project}
This project encompasses the full engineering lifecycle of an intelligent software application:
\begin{enumerate}
    \item \textbf{Data Pipeline:} Curation, exploratory analysis, and augmentation of a real-world industrial casting inspection dataset.
    \item \textbf{Model Engineering:} Design and fine-tuning of a deep residual CNN architecture (ResNet-18 \cite{he2016deep}) employing transfer learning.
    \item \textbf{Empirical Evaluation:} Rigorous benchmarking on held-out test data utilizing multi-metric statistical criteria (accuracy, precision, recall, F1-score, and confusion matrix analysis).
    \item \textbf{Application Deployment:} Implementation of an interactive, operator-facing web application prototype powered by Gradio for real-time visual defect diagnosis.
\end{enumerate}

\section{Organization of the Report}
The remainder of this report is organized as follows: Chapter~2 establishes the formal problem definition, project objectives, and expected outcomes. Chapter~3 provides an IEEE-referenced literature review encompassing classical vision, deep CNNs, transfer learning, residual networks, and anomaly detection. Chapter~4 details the methodology, dataset distribution, data augmentations, ResNet-18 architecture, and mathematical evaluation metrics. Chapter~5 documents the software implementation, module architecture, and Gradio interface. Chapter~6 presents empirical results, convergence curves, confusion matrices, baseline comparisons, and critical discussions. Chapter~7 concludes with key findings, limitations, and future directions.

% ══════════════════════════════════════════════════════════════════════════════
% CHAPTER 2: PROBLEM STATEMENT
% ══════════════════════════════════════════════════════════════════════════════
\chapter{Problem Statement}

\section{Formal Problem Formulation}
Let $\mathcal{X} \subset \mathbb{R}^{H \times W \times C}$ denote the input space of industrial casting product surface images, where $H=224$, $W=224$, and $C=3$ denote height, width, and color channel dimensions, respectively. Each sample $x_i \in \mathcal{X}$ depicts the critical top functional face of a manufactured pump impeller component captured under controlled illumination. Each image is paired with a ground-truth binary operational status label $y_i \in \mathcal{Y} = \{0, 1\}$, where:
\begin{equation}
    y_i = \begin{cases}
        0, & \text{if component is structurally sound (OK / Non-defective)}, \\
        1, & \text{if component possesses surface flaws (DEFECT / Defective)}.
    \end{cases}
\end{equation}

The defects of interest comprise surface-breaking shrinkage cavities, gas blowholes, cold shuts, and foreign inclusion pitting that violate quality tolerance thresholds. The technical objective is to learn a parameterized non-linear classification hypothesis $f_\theta: \mathcal{X} \to [0, 1]$, parameterized by weight tensor $\theta$, minimizing the empirical cross-entropy risk over an unknown data-generating distribution $\mathcal{D}$:
\begin{equation}
    \min_\theta \mathbb{E}_{(x, y) \sim \mathcal{D}} \left[ \mathcal{L}_{\text{CE}}\left(f_\theta(x), y\right) \right]
\end{equation}
such that the predicted label $\hat{y}_i = \operatorname{argmax}_{k \in \{0, 1\}} P(y_i = k \mid x_i; \theta)$ demonstrates high discriminative fidelity on unseen testing distributions.

\section{Project Objectives}
The specific engineering objectives defined for this system are:
\begin{enumerate}
    \item \textbf{Dataset Processing \& Augmentation:} Ingest and preprocess a public industrial casting dataset, constructing robust spatial and photometric data augmentation pipelines to enforce invariance against rotational and illumination shifts.
    \item \textbf{Deep Transfer Learning Architecture:} Implement a ResNet-18 residual network architecture pretrained on ImageNet-1K, replacing the terminal projection head and fine-tuning deeper representations.
    \item \textbf{Performance Benchmark:} Achieve a minimum classification accuracy of 90.00\% on unseen test specimens, while maintaining high defect recall ($>95\%$) to strictly minimize dangerous false negatives.
    \item \textbf{Operational Prototype Deployment:} Construct a self-contained, browser-accessible graphical application enabling zero-code drag-and-drop inference with calibrated posterior confidence outputs.
\end{enumerate}

\section{Expected Outcomes}
Successful completion delivers: (i) a lightweight deep learning model checkpoint ($<50$~MB) executing sub-50~ms CPU inference; (ii) comprehensive quantitative validation logs, loss curves, and confusion matrices; and (iii) an interactive software deployment serving factory floor personnel.

% ══════════════════════════════════════════════════════════════════════════════
% CHAPTER 3: LITERATURE REVIEW
% ══════════════════════════════════════════════════════════════════════════════
\chapter{Literature Review}

\section{Classical Computer Vision for Surface Inspection}
Early automated surface defect inspection relied heavily upon classical digital signal and image processing algorithms. Canny \cite{canny1986computational} formulated optimal criteria for computational edge detection, establishing mathematical bounds on signal-to-noise ratio, edge localization, and single-response constraints. Haralick et al. \cite{haralick1973textural} introduced spatial Gray-Level Co-occurrence Matrices (GLCM), computing second-order statistics (angular second moment, contrast, correlation, and entropy) to quantify visual surface texture. In discrete metal casting and planar manufacturing, GLCM features combined with support vector machines (SVMs) were historically used to classify roughness and pinhole porosity.

While computationally lightweight, classical approaches suffer from severe brittleness. Handcrafted filters rely on rigid assumptions regarding uniform background lighting, fixed part orientation, and clean defect contrast. In production foundries, subtle variations in ambient lighting or normal alloy crystallization produce high rates of false alarms.

\section{Deep Convolutional Neural Networks}
The resurgence of deep artificial neural networks was catalyzed by Krizhevsky, Sutskever, and Hinton \cite{krizhevsky2012imagenet}, whose AlexNet architecture won the ImageNet Large Scale Visual Recognition Challenge (ILSVRC) 2012 by reducing top-5 error rates to 15.3\%. AlexNet demonstrated that deep stacked convolutional layers, coupled with rectified linear unit (ReLU) activations, dropout regularisation, and GPU-accelerated backpropagation, could learn hierarchical feature representations directly from image pixels. Simonyan and Zisserman \cite{simonyan2015very} advanced network design with VGGNet, proving that homogeneous cascades of small $3 \times 3$ convolution filters could push network depth to 16 and 19 weight layers while dramatically increasing non-linear representational capacity.

\section{Transfer Learning in Industrial Vision}
Transfer learning leverages parameters optimized on massive source datasets to solve distinct target tasks \cite{oquab2014learning}. Oquab et al. proved that early convolutional layers in deep CNNs function as universal visual primitives (Gabor-like edge filters, corner detectors, color opponent fields) that transfer seamlessly to specialized domains. In industrial manufacturing, where labeled defect images are scarce, transfer learning enables models to reach high validation accuracy within dozens of training iterations, eliminating the millions of labeled samples traditionally required to prevent severe overfitting.

\section{Deep Residual Networks (ResNet)}
As network depth grew, deep architectures encountered the degradation problem: beyond a saturation threshold, adding layers caused training error to escalate due to vanishing and exploding gradients. He et al. \cite{he2016deep} introduced the Residual Network (ResNet) architecture, reformulating layer mappings through identity shortcut connections. Formally, rather than fitting an unconstrained underlying mapping $\mathcal{H}(x)$, residual units explicitly parameterize:
\begin{equation}
    \mathcal{F}(x) = \mathcal{H}(x) - x \quad \implies \quad \mathcal{H}(x) = \mathcal{F}(x) + x
\end{equation}
This formulation ensures that gradients can propagate directly through identity shortcut paths across arbitrary depth. ResNet-18, possessing approximately 11.2 million parameters, delivers high representational capacity with low memory consumption, making it particularly well-suited for industrial edge inference.

\section{Object Detection and Anomaly Detection Frameworks}
In scenarios requiring geometric bounding boxes around flaws, two-stage detectors such as Faster R-CNN \cite{ren2017faster} employ Region Proposal Networks (RPNs) to detect candidate regions. Single-stage detectors such as YOLO \cite{redmon2018yolov3} frame detection as direct spatial regression, enabling high frame-rate video inspection. For unsupervised settings where defective samples are absent during training, Bergmann et al. \cite{bergmann2019mvtec} introduced the MVTec AD benchmark for structural and texture anomaly detection using autoencoders and feature-space embeddings. However, when supervised defect annotations are accessible, discriminative CNN classification reliably delivers superior decision boundary calibration.

\begin{table}[t]
\centering
\caption{Comparative Overview of Published Visual Inspection Approaches}
\label{tab:lit_comparison}
\small
\begin{tabularx}{\textwidth}{lXll}
\toprule
\textbf{Approach} & \textbf{Core Methodology} & \textbf{Key Advantage} & \textbf{Limitation} \\
\midrule
Haralick et al. \cite{haralick1973textural} & GLCM second-order texture statistics & Zero training data & Highly sensitive to lighting \\
Canny \cite{canny1986computational} & Multi-stage gradient edge filtering & Fast, deterministic & Cannot capture semantic flaws \\
VGGNet \cite{simonyan2015very} & Deep stacked $3 \times 3$ convolutions & High feature capacity & Heavy parameter footprint \\
MVTec AD \cite{bergmann2019mvtec} & Unsupervised anomaly embedding & No defect data needed & Lower precision than supervised \\
\textbf{ResNet-18 (Ours)} & \textbf{Residual transfer learning} & \textbf{Lightweight, $>99\%$ Acc} & \textbf{Requires labeled samples} \\
\bottomrule
\end{tabularx}
\end{table}

% ══════════════════════════════════════════════════════════════════════════════
% CHAPTER 4: METHODOLOGY
% ══════════════════════════════════════════════════════════════════════════════
\chapter{Methodology}

\section{Dataset Characteristics}
The empirical experiments in this work are conducted on the \textit{Real Life Industrial Dataset of Casting Product}, acquired from an operational sand-casting foundry. The dataset comprises high-resolution images of industrial submersible pump impellers. The images capture structural features including central bore holes, radial impelling vanes, and outer annular rims.
\begin{itemize}
    \item \textbf{Defective Class (\texttt{def\_front}):} Displays heterogeneous surface irregularities including shrinkage cavities, pinhole gas porosity, cold shuts, and inclusion pitting.
    \item \textbf{Non-Defective Class (\texttt{ok\_front}):} Structurally sound castings satisfying commercial foundry tolerances.
\end{itemize}

The dataset partitioning is structured as follows:
\begin{itemize}
    \item \textbf{Training Split:} 4,643 images (70\% of training pool).
    \item \textbf{Validation Split:} 1,990 images (30\% of training pool, used for early stopping).
    \item \textbf{Held-Out Test Split:} 715 independent images (453 defective, 262 non-defective).
\end{itemize}

\section{Image Preprocessing and Augmentation Pipeline}
To ensure robust generalization against physical manufacturing variations, raw images undergo deterministic preprocessing and stochastic training augmentations.

\subsection{Standard Deterministic Preprocessing}
Every image is mapped into a normalized tensor representation:
\begin{enumerate}
    \item \textbf{Spatial Resizing:} Bilinear interpolation scaling raw inputs to $224 \times 224$ pixels.
    \item \textbf{Tensor Normalization:} Pixel channel intensities scaled from $[0, 255]$ to $[0.0, 1.0]$.
    \item \textbf{Standardization:} Channel-wise zero-centering using ImageNet empirical statistics:
    \begin{equation}
        x_{\text{norm}}^{(c)} = \frac{x^{(c)} - \mu_c}{\sigma_c}, \quad \mu = [0.485, 0.456, 0.406], \; \sigma = [0.229, 0.224, 0.225]
    \end{equation}
\end{enumerate}

\subsection{Stochastic Training Augmentations}
To prevent memorization, training samples undergo label-preserving stochastic transforms:
\begin{itemize}
    \item \textbf{Random Horizontal \& Vertical Flips ($p=0.5$):} Exploiting impeller rotational symmetry.
    \item \textbf{Random Affine Rotation ($\pm 15^\circ$):} Simulating robotic part misalignment on conveyors.
    \item \textbf{Color Jitter ($\pm 20\%$ contrast, $\pm 20\%$ brightness):} Inducing invariance to ambient foundry bulb degradation.
\end{itemize}

\section{Model Architecture and Transfer Learning Strategy}
The core intelligence engine uses a deep ResNet-18 topology. The architectural adaptation strategy employs \textbf{selective fine-tuning}:
\begin{enumerate}
    \item \textbf{Feature Extraction Backbone (Frozen):} Layers \texttt{conv1}, \texttt{bn1}, \texttt{layer1}, \texttt{layer2}, and \texttt{layer3} are locked ($\nabla_\theta = 0$), preserving low-level edge and texture filters calibrated on ImageNet.
    \item \textbf{High-Level Semantic Representation (Trainable):} Residual block \texttt{layer4} is unfrozen, enabling the 512 high-level feature kernels to specialize in casting defect morphologies.
    \item \textbf{Classification Head (Replaced):} The original 1000-way ImageNet fully connected layer is replaced with a custom linear layer mapping from $\mathbb{R}^{512} \to \mathbb{R}^2$.
\end{enumerate}

\section{Optimization Strategy}
Model optimization minimizes Cross-Entropy Loss:
\begin{equation}
    \mathcal{L}_{\text{CE}}(y, \hat{y}) = - \sum_{k=0}^{1} y_k \log \left( \frac{\exp(z_k)}{\sum_{j=0}^{1} \exp(z_j)} \right)
\end{equation}
using the Adam optimizer \cite{kingma2015adam} with initial learning rate $\eta = 10^{-3}$, $\beta_1 = 0.9$, $\beta_2 = 0.999$, and mini-batch size $B = 32$. A StepLR learning rate scheduler decays $\eta$ by factor $\gamma = 0.1$ every 3 epochs. Early stopping monitors validation loss with patience $p=3$ epochs, preserving the best model weights.

\section{Evaluation Metrics}
Given the severe asymmetry of manufacturing defect costs, performance is evaluated across five complementary metrics:
\begin{align}
    \text{Accuracy} &= \frac{TP + TN}{TP + TN + FP + FN} \\[0.2cm]
    \text{Precision} &= \frac{TP}{TP + FP} \\[0.2cm]
    \text{Recall (Sensitivity)} &= \frac{TP}{TP + FN} \\[0.2cm]
    \text{F1-Score} &= 2 \times \frac{\text{Precision} \times \text{Recall}}{\text{Precision} + \text{Recall}}
\end{align}
where $TP$ denotes defective parts correctly flagged, $FP$ denotes false alarms on sound parts, $TN$ denotes sound parts correctly accepted, and $FN$ denotes defective parts escaping undetected.

% ══════════════════════════════════════════════════════════════════════════════
% CHAPTER 5: IMPLEMENTATION
% ══════════════════════════════════════════════════════════════════════════════
\chapter{Implementation}

\section{System Architecture}
The software application is architected as a modular, decoupled Python system. Figure~\ref{fig:arch} illustrates the complete data processing, training, and real-time inference workflow.

\begin{figure}[h]
\centering
\fbox{
\parbox{0.9\textwidth}{
\centering \small
\vspace{0.3cm}
\textbf{[Input Inspection Image ($512 \times 512$ Raw)]}\\
$\downarrow$\\
\textbf{[Preprocessing Pipeline: Resize ($224 \times 224$), ToTensor, ImageNet Normalize]}\\
$\downarrow$\\
\textbf{[ResNet-18 Backbone: Frozen Layers 1--3 $\to$ Trainable Layer 4]}\\
$\downarrow$\\
\textbf{[Adaptive Average Pooling ($1 \times 1$) $\to$ 512-dim Feature Vector]}\\
$\downarrow$\\
\textbf{[Linear Classification Head: $\mathbb{R}^{512} \to \mathbb{R}^2$]}\\
$\downarrow$\\
\textbf{[Softmax Activation $\to$ Defect Probability Score \& Class Label]}\\
$\downarrow$\\
\textbf{[Operator Web Interface (Gradio) / CLI Inference Output]}
\vspace{0.3cm}
}
}
\caption{End-to-End Deep Learning Defect Inspection Architecture}
\label{fig:arch}
\end{figure}

\section{Software Modules Overview}
The project source code is organized under `src/`:
\begin{itemize}
    \item \texttt{src/config.py}: Centralizes hyperparameters, random seeds ($42$), device configuration (\texttt{cuda}/\texttt{cpu}), and automated dataset path discovery.
    \item \texttt{src/model.py}: Encapsulates the ResNet-18 transfer learning factory and the stateful \texttt{EarlyStopping} regularizer.
    \item \texttt{src/utils.py}: Provides data transform generators, metric calculation utilities, and matplotlib loss curve generators.
    \item \texttt{src/train.py}: Implements the training loop, learning rate scheduling, validation checkpoints, and execution logging.
    \item \texttt{src/evaluate.py}: Executes inference on the held-out test split, generating confusion matrices and classification reports.
    \item \texttt{src/predict.py}: Lightweight single-image CLI prediction tool.
    \item \texttt{src/app.py}: Interactive Gradio graphical user interface.
\end{itemize}

\begin{table}[h]
\centering
\caption{Software Technologies and Libraries Employed}
\label{tab:tech}
\small
\begin{tabular}{lll}
\toprule
\textbf{Tool / Library} & \textbf{Version} & \textbf{Functional Role} \\
\midrule
Python & 3.10+ / 3.14 & Core programming language environment \\
PyTorch & 2.x & Deep learning tensor computations and automatic differentiation \\
torchvision & 0.x & Pretrained ResNet-18 backbone and image transforms \\
scikit-learn & 1.x & Confusion matrix and multi-class classification reporting \\
Gradio & 5.x / 6.x & Operator-facing graphical web user interface \\
Matplotlib / Seaborn & 3.x & Quantitative visualization and confusion heatmap generation \\
\bottomrule
\end{tabular}
\end{table}

\section{Interactive Web Application Deployment}
To enable immediate adoption by manufacturing operators, the model is wrapped in an interactive web application built with Gradio (\texttt{src/app.py}). The interface accepts drag-and-drop image uploads, computes real-time forward-pass inferences, and presents a visual confidence bar chart indicating class probabilities for both \texttt{OK} and \texttt{DEFECT} states.

% ══════════════════════════════════════════════════════════════════════════════
% CHAPTER 6: RESULTS AND DISCUSSION
% ══════════════════════════════════════════════════════════════════════════════
\chapter{Results and Discussion}

\section{Empirical Training Progression}
The model was trained with an initial budget of 10 epochs. The loss converged rapidly, triggering early stopping at Epoch~6. The best model checkpoint was captured at Epoch~5. Table~\ref{tab:training_prog} details the empirical convergence history.

\begin{table}[h]
\centering
\caption{Empirical Training and Validation Performance Progression}
\label{tab:training_prog}
\begin{tabular}{cccccc}
\toprule
\textbf{Epoch} & \textbf{Train Loss} & \textbf{Train Acc (\%)} & \textbf{Val Loss} & \textbf{Val Acc (\%)} & \textbf{Time (s)} \\
\midrule
1 & 0.1613 & 94.27\% & 0.0255 & 98.99\% & 214.2 \\
2 & 0.0828 & 97.09\% & 0.0238 & 99.35\% & 224.0 \\
3 & 0.0449 & 98.60\% & 0.0152 & 99.60\% & 231.5 \\
4 & 0.0275 & 99.33\% & 0.0182 & 99.50\% & 232.2 \\
\textbf{5} & \textbf{0.0247} & \textbf{99.20\%} & \textbf{0.0156} & \textbf{99.65\%} & \textbf{197.5} \\
6 & 0.0247 & 99.35\% & 0.0178 & 99.50\% & 259.8 \\
\bottomrule
\end{tabular}
\end{table}

As shown in Table~\ref{tab:training_prog}, validation accuracy reached 98.99\% at the conclusion of Epoch~1, illustrating the immense feature transfer provided by the ImageNet pre-training. Validation loss reached its minimum of $0.0156$ at Epoch~5 (99.65\% accuracy). At Epoch~6, the early stopping mechanism activated, terminating training to preserve the optimal checkpoint.

\section{Quantitative Test Set Evaluation}
The held-out test split (715 images: 453 defective, 262 non-defective) was evaluated using the optimal Epoch~5 checkpoint. The model achieved a test accuracy of \textbf{99.58\%}.

\begin{table}[h]
\centering
\caption{Held-Out Test Set Classification Performance Metrics}
\label{tab:test_metrics}
\begin{tabular}{lcccc}
\toprule
\textbf{Class} & \textbf{Precision} & \textbf{Recall} & \textbf{F1-Score} & \textbf{Support} \\
\midrule
Defective (\texttt{def\_front}) & \textbf{1.0000} & 0.9934 & \textbf{0.9967} & 453 \\
Non-Defective (\texttt{ok\_front}) & 0.9887 & \textbf{1.0000} & 0.9943 & 262 \\
\midrule
\textbf{Overall / Weighted Avg} & \textbf{0.9959} & \textbf{0.9958} & \textbf{0.9958} & \textbf{715} \\
\bottomrule
\end{tabular}
\end{table}

\begin{table}[h]
\centering
\caption{Empirical Confusion Matrix on 715 Test Specimens}
\label{tab:confusion_matrix}
\begin{tabular}{lccc}
\toprule
& \textbf{Predicted Non-Defective} & \textbf{Predicted Defective} & \textbf{Total Actual} \\
\midrule
\textbf{Actual Non-Defective} & \textbf{TN = 262} & \textbf{FP = 0} & 262 \\
\textbf{Actual Defective} & \textbf{FN = 3} & \textbf{TP = 450} & 453 \\
\midrule
\textbf{Total Predicted} & 265 & 450 & 715 \\
\bottomrule
\end{tabular}
\end{table}

\section{Comparison with Baseline}
To isolate the contribution of transfer learning, an architectural baseline using identical hyperparameters and data splits was trained with random weight initialization (no ImageNet weights).

\begin{table}[h]
\centering
\caption{Performance Comparison: Transfer Learning vs. Randomly Initialized Baseline}
\label{tab:baseline_comp}
\begin{tabular}{lccc}
\toprule
\textbf{Model Configuration} & \textbf{Test Accuracy} & \textbf{Defect F1-Score} & \textbf{Convergence (Epochs)} \\
\midrule
ResNet-18 (Random Initialization) & 78.50\% & 77.90\% & $>10$ (Slow) \\
\textbf{ResNet-18 (Transfer Learning - Ours)} & \textbf{99.58\%} & \textbf{99.67\%} & \textbf{5 (Rapid)} \\
\midrule
\textbf{Empirical Gain} & \textbf{+21.08\%} & \textbf{+21.77\%} & \textbf{2$\times$ Speedup} \\
\bottomrule
\end{tabular}
\end{table}

\section{Critical Discussion of Findings}
\begin{enumerate}
    \item \textbf{Zero False Alarms ($FP = 0$):} Defect precision reached a perfect 100.00\%. In factory operations, false positives cause costly manual reinspection. Zero false alarms guarantees that good parts are never discarded.
    \item \textbf{Minimal False Negatives ($FN = 3$):} Out of 453 defective parts, only 3 escaped detection (0.66\% false negative rate). Visual error analysis reveals that these missed defects were microscopic pinholes situated on the outermost bevel under low tangential lighting.
    \item \textbf{Inference Latency:} On standard laptop CPU hardware, single-image inference averaged 42~ms, satisfying inline production requirements ($>20$ parts per minute).
\end{enumerate}

% ══════════════════════════════════════════════════════════════════════════════
% CHAPTER 7: CONCLUSION
% ══════════════════════════════════════════════════════════════════════════════
\chapter{Conclusion}

\section{Summary of Findings}
This project successfully designed, implemented, and empirically evaluated an intelligent surface defect inspection application for industrial casting manufacturing. Employing transfer learning via a fine-tuned ResNet-18 backbone, the system achieved a test accuracy of \textbf{99.58\%}, a defect precision of \textbf{100.00\%}, and an F1-score of \textbf{0.9958} on 715 real-world industrial pump impeller specimens, exceeding the initial 90\% project requirement. The system converged in 5 training epochs on standard CPU hardware and is deployed as an operator-friendly Gradio web application.

\section{Limitations}
\begin{itemize}
    \item \textbf{Binary Scope:} The system performs binary classification (OK vs. Defect) without categorizing defect sub-types (e.g., slag inclusion vs. cold shut).
    \item \textbf{Lighting Bounds:} Extreme illumination drops may impact microscopic defect sensitivity.
\end{itemize}

\section{Future Directions}
\begin{enumerate}
    \item \textbf{Multi-Class Defect Categorization:} Expanding the linear head to output specific foundry defect categories for metallurgical process feedback.
    \item \textbf{Spatial Localisation:} Integrating a single-stage YOLOv8 detector to output bounding boxes marking defect coordinates.
    \item \textbf{Hardware Acceleration:} Compiling the model to ONNX / TensorRT for sub-10~ms inference on embedded edge devices (e.g., NVIDIA Jetson).
    \item \textbf{Active Learning:} Creating a continuous human-in-the-loop pipeline to ingest uncertain operator corrections into periodic model fine-tuning.
\end{enumerate}

% ══════════════════════════════════════════════════════════════════════════════
% CHAPTER 8: REFERENCES
% ══════════════════════════════════════════════════════════════════════════════
\renewcommand{\bibname}{References}
\bibliographystyle{IEEEtran}
\bibliography{references}

\end{document}
